\subsection{What the reconstructed ring says about homotopy}\label{sec:homotopy-recovery}
The reconstruction theorem identifies the algebraic input to existing
integral homotopy-recovery theorems. Here we state exactly the consequences
and distinguish the two categories in which derived maps are taken.

\begin{proof}[Proof of Corollary~\ref{cor:homotopy-recovery}]
For each $n$, the function ring $A_X^n=\operatorname{Map}(X_n,\Z)$
is torsion-free and has pointwise binomial operations
$f\mapsto\binom f r$. These operations are uniquely determined by
\[
 r!\binom f r=f(f-1)\cdots(f-r+1),
\]
because multiplication by $r!$ is injective. Every ring map between
such rings preserves them. Consequently the cosimplicial ring of
Theorem~\ref{thm:reconstruction} canonically determines its
cosimplicial binomial-ring structure.

Horel's model category of cosimplicial binomial rings has weak
equivalences the quasi-isomorphisms of normalized cochain complexes
\cite[Theorem~3.7, p.~7]{Horel}. His adjunction sends a simplicial
set $X$ to $\Z^X$ and has derived right adjoint
\[
 \langle A\rangle=\mathbb R\Map_{\mathsf{cBRing}}(A,\Z)
 \qquad\text{\cite[Section~4, p.~9]{Horel}}.
\]
For connected nilpotent finite-type $Y$, his Theorem~5.4, p.~13,
says that the unit $Y\to\langle A_Y\rangle$ is a weak equivalence.
The derived adjunction therefore gives
\[
 \mathbb R\Map_{\mathsf{cBRing}}(A_Y,A_X)
 \simeq\mathbb R\Map_{\mathsf{Spaces}}(X,\langle A_Y\rangle)
 \simeq\mathbb R\Map_{\mathsf{Spaces}}(X,Y).
\]
These are applications of Horel's theorem, rather than a new
homotopy-recovery theorem.

For the integral $E_\infty$ cochain category, Mandell's Main Theorem
\cite[p.~213]{Mandell} says that the cochain algebra detects weak
homotopy types of finite-type nilpotent spaces. His Theorem~0.1,
p.~214, gives faithfulness. His Theorem~0.2, p.~215, describes the
extra maps: the right adjoint applied to $C^*(Y;\Z)$ is the homotopy
pullback of the free loop space of the finite completion $Y^\wedge$
along $Y\to Y^\wedge$. Thus for simply connected $Y$ its $n$th
homotopy group is
\[
 \pi_nY\oplus\widehat{\pi_{n+1}Y}.
\]
In particular, for $X=S^2$ and $Y=S^3$, the homotopy set $[X,Y]$
is a singleton, whereas the set of derived $E_\infty$ cochain maps
is $\widehat\Z$ and is uncountable; this is Mandell's explicit
example on p.~216. The two mapping-space assertions refer to
different derived categories.

Finally let $X$ have finitely many components, each nilpotent and
of finite type. The primitive idempotents of
$H^0(A_X)=\Z^{\pi_0X}$ recover the component set. The normalized
complex is nonnegative, so a degree-zero cohomology class has a
unique cocycle representative. For each primitive idempotent this
representative is the characteristic function of one component on
vertices. Its cocycle condition makes it constant on the endpoints
of every edge. Consequently all its cosimplicial images in degree
$n$ agree on the vertices of each simplex and give the characteristic
idempotent of that component of $X_n$. These idempotents split
$A_X$ canonically, up to permutation, into the function rings of its
components. Apply the connected object-detection theorem to each
factor. The zero function ring distinguishes the empty space.
\end{proof}


To\"en's integral affinization of a simply connected finite-type
space gives another comparison. Its homotopy sheaves are the homotopy
groups tensored with the Hilbert additive group scheme, and its
integral points have a natural retraction onto the original space
\cite[Theorem~4.1, p.~11; Corollaries~5.4--5.5, pp.~16--17]{Toen}.
The unit to those points need not be an equivalence: Corollary~5.1,
p.~14, gives additional completed homotopy groups, including
$\widehat\Z$ in degree two for $S^3$. The binomial category in
Corollary~\ref{cor:homotopy-recovery} has the stated unit equivalence.
